../cictikz/src/cictikz/data/tex/cictikz_preamble.tex

% Differential ring oscillator drawn with OTA symbols: two differential
% delay stages in a loop, the loop inversion is a wire swap between
% stage one and stage two, and a third stage buffers the output.

\begin{circuitikz}[american, thick]

  % differential stage as an OTA triangle: inputs left (+ top, - bottom),
  % outputs right (- top, + bottom)
  \newcommand{\difftri}[3]{%
    \draw ({#1-0.9},{#2-1.0}) -- ({#1-0.9},{#2+1.0}) -- ({#1+0.9},#2) -- cycle;
    \coordinate (#3inp) at ({#1-0.9},{#2+0.45});
    \coordinate (#3inn) at ({#1-0.9},{#2-0.45});
    \coordinate (#3outn) at ({#1+0.09},{#2+0.45});
    \coordinate (#3outp) at ({#1+0.09},{#2-0.45});
    \node[anchor=west, scale=0.8] at ({#1-0.85},{#2+0.45}) {$+$};
    \node[anchor=west, scale=0.8] at ({#1-0.85},{#2-0.45}) {$-$};
    \node[anchor=east, scale=0.8] at ({#1+0.05},{#2+0.45}) {$-$};
    \node[anchor=east, scale=0.8] at ({#1+0.05},{#2-0.45}) {$+$};}

  \difftri{0}{0}{a}
  \difftri{4.4}{0}{b}
  \difftri{8.8}{0}{c}

  % stage one to stage two: the wires cross, and that swap is the -1
  % that makes the loop oscillate
  \draw (aoutn) -- (1.6,0.45) -- (3.0,-0.45) -- (binn);
  \draw (aoutp) -- (1.6,-0.45) -- (3.0,0.45) -- (binp);

  % stage two outputs, tapped by the buffer stage
  \draw (boutn) -- ++(1.3,0) coordinate (w3);
  \draw (boutp) -- ++(1.3,0) coordinate (w4);
  \fill (w3) circle (0.06);
  \fill (w4) circle (0.06);
  \draw (w3) -- (cinp);
  \draw (w4) -- (cinn);

  % feedback: stage two straight back to stage one
  \draw (w3) -- (w3 |- 0,2.0) -- (-2.2,2.0) -- (-2.2,0.45) -- (ainp);
  \draw (w4) -- (w4 |- 0,-2.0) -- (-2.2,-2.0) -- (-2.2,-0.45) -- (ainn);

  % buffered output
  \draw (coutn) -- ++(1.0,0) coordinate (q1);
  \draw (q1) to [short,-o] ++(0.3,0) node[anchor=west] {$v_{o-}$};
  \draw (coutp) -- ++(1.0,0) coordinate (q2);
  \draw (q2) to [short,-o] ++(0.3,0) node[anchor=west] {$v_{o+}$};

\end{circuitikz}
\end{document}
